\documentclass[10pt, aspectratio=169]{beamer}
\input{../../_en_preambulo_beamer}
% Comandos y macros estadísticas compartidas para las presentaciones Beamer en inglés (Metropolis)

% Conjuntos numéricos básicos
\newcommand{\R}{\mathbb{R}}
\newcommand{\N}{\mathbb{N}}
\newcommand{\Q}{\mathbb{Q}}
\newcommand{\Z}{\mathbb{Z}}
\newcommand{\set}[1]{\left\{ #1 \right\}}
\newcommand{\abs}[1]{\left|#1\right|}
\newcommand{\norm}[1]{\left\| #1 \right\|}

% Comandos estadísticos y probabilísticos del curso
\newcommand{\Var}{\operatorname{Var}}
\newcommand{\cov}{\operatorname{Cov}}
\newcommand{\corr}{\rho}
\newcommand{\comb}[2]{\begin{pmatrix} #1 \\ #2 \end{pmatrix}}
\newcommand{\s}{\sigma}
\newcommand{\particion}{\mathfrak{P}}
\newcommand{\card}[1]{\#\left( #1 \right)}
\newcommand{\seq}[3]{\set{#1}_{#2}^{#3}}
\newcommand{\E}{\mathbb{E}}
\newcommand{\Prob}{\mathbb{P}}

% Entornos pedagógicos y cajas en inglés académico natural (soportando tanto nombres en inglés como en español)
\newenvironment{discussion}{%
  \begin{alertblock}{Discussion Question}%
}{%
  \end{alertblock}%
}
\newenvironment{discusion}{%
  \begin{alertblock}{Discussion Question}%
}{%
  \end{alertblock}%
}

\newenvironment{reflection}{%
  \begin{alertblock}{Classroom Reflection}%
}{%
  \end{alertblock}%
}
\newenvironment{reflexion}{%
  \begin{alertblock}{Classroom Reflection}%
}{%
  \end{alertblock}%
}

\newenvironment{paradox}{%
  \begin{alertblock}{Exploratory Paradox}%
}{%
  \end{alertblock}%
}
\newenvironment{paradoja}{%
  \begin{alertblock}{Exploratory Paradox}%
}{%
  \end{alertblock}%
}

\newenvironment{motivation}{%
  \begin{exampleblock}{Motivation: Data Science in Practice}%
}{%
  \end{exampleblock}%
}
\newenvironment{motivacion}{%
  \begin{exampleblock}{Motivation: Data Science in Practice}%
}{%
  \end{exampleblock}%
}

\newenvironment{quickchallenge}{%
  \begin{block}{Quick Team Challenge}%
}{%
  \end{block}%
}
\newenvironment{retoclase}{%
  \begin{block}{Quick Team Challenge}%
}{%
  \end{block}%
}


\title[Implementation in Python]{Section 09.06: Implementation in Python with \texttt{statsmodels}}
\subtitle{From Manual Formulas to the Professional Summary Table}
\author[J. Castillo Colmenares]{Juliho Castillo Colmenares}
\institute{Tecnológico de Monterrey}
\date{\vspace{-1.5cm}}

\begin{document}

% Slide 1: Title
\begin{frame}[plain]
  \titlepage
\end{frame}

% Slide 2: Roadmap
\begin{frame}{Roadmap --- Chapter 09: Linear and Multiple Regressions}
  \scriptsize
  Where are we in the modular development of this chapter?
  \begin{itemize}\itemsep=0.02em
    \item \textbf{09.01} Correlation as the Premise of Regression
    \item \textbf{09.02} Introduction to Linear Regression
    \item \textbf{09.03} The Mathematics of Regression: OLS Derivation and $R^2$
    \item \textbf{09.04} Linear Regression on Simulated Data
    \item \textbf{09.05} Optimal Coefficients, $t$/$F$ Tests, and RSE
    \item \textbf{09.06} Implementation in Python with \texttt{statsmodels} \emph{(Today)}
    \item \textbf{09.07} Multiple Regression, the Normal Equation, and Ridge/Lasso
    \item \textbf{09.08} Model Validation and $k$-fold Cross-Validation
    \item \textbf{09.09} Residual Diagnostics and Multicollinearity (VIF)
    \item \textbf{09.10} Regression with \texttt{scikit-learn}
    \item \textbf{09.11} Categorical Variables and Dummy Variables
    \item \textbf{09.12} Nonlinear Transformations and Polynomial Regression
  \end{itemize}
  \pause
  \vspace{-0.05cm}
  \begin{block}{Session Objective}
    Recognize how the professional \texttt{statsmodels} library condenses into a single summary table exactly the same calculations manually derived in Sections 09.03-09.05: coefficients, standard errors, $t$/$F$ statistics, $R^2$, and RSE.
  \end{block}
\end{frame}

% Slide 3: Motivation
\begin{frame}{From Manual Formulas to the Professional Tool}
  \footnotesize
  \begin{motivacion}
    We have manually derived and computed every component of a regression model. In professional practice, libraries like \texttt{statsmodels} automate these calculations in a single line of code, producing a standardized summary table used across industry and academia.
  \end{motivacion}

  \vspace{0.15cm}
  \pause
  Python offers two main routes for regression: \texttt{statsmodels.formula.api} (statistical/inferential focus, similar to R) and \texttt{scikit-learn} (machine-learning focus, Section 09.10).
\end{frame}

% Slide 4: The ols method
\begin{frame}[fragile]{The \texttt{statsmodels} \texttt{ols} Method}
  \footnotesize
  \begin{block}{Fitting with Formula Notation}
    \begin{verbatim}
import statsmodels.formula.api as smf
model = smf.ols(formula='Sales~TV', data=advert).fit()
print(model.params)
    \end{verbatim}
  \end{block}
  \pause

  \vspace{0.1cm}
  \begin{alertblock}{Output}
    \texttt{Intercept} $\approx7.03$, \texttt{TV} $\approx0.0475$ --- exactly the same $\hat\alpha,\hat\beta$ we would compute by hand with the Section 09.03 formulas.
  \end{alertblock}
\end{frame}

% Slide 5: The full summary table
\begin{frame}{The \texttt{model.summary()} Table}
  \footnotesize
  A single call (\texttt{model.summary()}) simultaneously reports:
  \begin{itemize}\itemsep=0.08cm
    \item \textbf{Coefficients} and their \textbf{standard errors} (Section 09.05). \pause
    \item \textbf{$t$-statistics and $p$-values} per coefficient (Section 09.05). \pause
    \item \textbf{$R^2$ and adjusted $R^2$} (Section 09.03). \pause
    \item Overall \textbf{$F$-statistic} and its $p$-value (Section 09.05). \pause
    \item \textbf{Durbin-Watson}, AIC, BIC --- diagnostics we will develop further in Section 09.09.
  \end{itemize}
\end{frame}

% Slide 6: Prediction and RSE
\begin{frame}[fragile]{Prediction and Residual Standard Error}
  \footnotesize
  \begin{block}{Predicting on New Data}
    \begin{verbatim}
sales_pred = model.predict(pd.DataFrame(advert['TV']))
    \end{verbatim}
  \end{block}
  \pause

  \vspace{0.1cm}
  \begin{alertblock}{RSE and Relative Error}
    With $\text{RSE}\approx3.25$ over an average sale of $\approx14.02$, the model's relative error is $\approx23\%$: substantial, and it motivates adding more predictors (\texttt{Radio}, \texttt{Newspaper}) --- the subject of Section 09.07.
  \end{alertblock}
\end{frame}

% Slide 7: Python Lab (Block 1)
\begin{frame}[fragile]{Python Lab: Reconstructing the Dataset}
  \scriptsize
  Synthetic reconstruction of the classic advertising TV/Sales example (\texttt{09.06\_statsmodels\_style\_summary.py}):
  {\lstset{basicstyle=\fontsize{3.6pt}{4.3pt}\selectfont\ttfamily,aboveskip=0.05em,belowskip=0.05em}
  \lstinputlisting[firstline=18, lastline=25]{../../code/09_regresiones/09.06_statsmodels_style_summary.py}}
\end{frame}

% Slide 8: Python Lab (Block 2)
\begin{frame}[fragile]{Python Lab: Full Summary Table from Scratch}
  \scriptsize
  Manually reproducing every entry of a \texttt{statsmodels} table using only \texttt{numpy}/\texttt{scipy}:
  {\lstset{basicstyle=\fontsize{3.0pt}{3.7pt}\selectfont\ttfamily,aboveskip=0.05em,belowskip=0.05em}
  \lstinputlisting[firstline=29, lastline=48]{../../code/09_regresiones/09.06_statsmodels_style_summary.py}}
\end{frame}

% Slide 9: Python Lab (Block 3)
\begin{frame}[fragile]{Python Lab: $R^2$, $F$-Statistic, and RSE}
  \scriptsize
  Computing the final global fit-quality metrics:
  {\lstset{basicstyle=\fontsize{3.4pt}{4.1pt}\selectfont\ttfamily,aboveskip=0.05em,belowskip=0.05em}
  \lstinputlisting[firstline=51, lastline=64]{../../code/09_regresiones/09.06_statsmodels_style_summary.py}}
\end{frame}

% Slide 10: Terminal Output
\begin{frame}[fragile]{Python Lab: Terminal Output}
  \scriptsize
  Standard output produced when running the lab script:
  \vspace{0.02cm}
  \begin{block}{Standard Console Output}
    \fontsize{3.4pt}{4.1pt}\selectfont
    \begin{verbatim}
=== Block 1: Reconstructing Advertising-Style Dataset ===
n=200 observations; mean(TV)=146.74, mean(Sales)=14.05

=== Block 2: Full OLS Summary Table ===
Coefficient    Estimate     Std Err         t       P>|t|
Intercept        6.0202      0.4670     12.8914    5.24e-28
TV               0.0547      0.0028     19.7788    9.01e-49

=== Block 3: R^2, F-Statistic, and RSE ===
R-squared = 0.6640
F-statistic = 391.2003 (p-value=9.01e-49)
RSE = 3.2595, mean(Sales) = 14.0533, relative error = 0.2319 (23.2%)
    \end{verbatim}
  \end{block}
\end{frame}

% Slide 11: Comparison with the book's example
\begin{frame}{Comparison: Our Engine vs. \texttt{statsmodels}}
  \footnotesize
  \begin{itemize}\itemsep=0.12cm
    \item Our synthetic reconstruction produces, for this particular sample, intercept $\approx6.02$, \texttt{TV} slope $\approx0.055$, $R^2\approx0.664$, RSE$\approx3.26$ with relative error $\approx23\%$ --- the same order of magnitude as the book's classic example (intercept $\approx7.03$, slope $\approx0.0475$, $R^2\approx0.61$); the specific differences are due to the sampling variability inherent to a distinct synthetic reconstruction of the original dataset. \pause
    \item This confirms that our manual OLS engine (Sections 09.03-09.05) computes \textbf{exactly} what \texttt{statsmodels} would report on the same data --- the library uses no different mathematics, it just automates the calculation. \pause
    \item The $\approx23\%$ relative error is substantial: it motivates incorporating more predictors into a multiple regression model (Section 09.07).
  \end{itemize}
\end{frame}

% Slide 12: Closing
\begin{frame}{Section Closing: Implementation with \texttt{statsmodels}}
  \begin{reflexion}
    All the machinery manually derived in Sections 09.03-09.05 --- OLS coefficients, $t$/$F$ tests, RSE --- is exactly what a professional library automates into a single summary table. Understanding the underlying mathematics is what allows one to correctly interpret that table, rather than treating it as a black box.
  \end{reflexion}

  \vspace{0.2cm}
  \pause
  \begin{block}{Key Takeaways}
    \begin{itemize}\itemsep=0.08cm
      \item \textbf{\texttt{statsmodels.formula.api.ols}:} R-style formula syntax, inferential focus.
      \item \textbf{\texttt{model.summary()}:} condenses coefficients, standard errors, $t$, $F$, $R^2$, and diagnostics into one table.
      \item \textbf{No additional "trick":} the library reproduces exactly the already-derived formulas, introducing no new mathematics.
    \end{itemize}
  \end{block}
\end{frame}

% Slide 13: Modular Perspective
\begin{frame}{Modular Perspective --- The Next Challenge}
  \scriptsize
  \begin{retoclase}
    Our single-predictor model's $\approx23\%$ relative error suggests that \texttt{TV} does not tell the whole story. Section 09.07 generalizes the model to \textbf{multiple} simultaneous predictors (\texttt{TV}, \texttt{Radio}, \texttt{Newspaper}) via the Normal Equation in matrix notation, and presents Ridge/Lasso for when those predictors are correlated with each other.
  \end{retoclase}

  \vspace{0.08cm}
  \pause
  \begin{alertblock}{Recommended Preparation}
    \begin{itemize}\itemsep=0.06cm
      \item Review basic matrix notation: vectors, matrices, transpose, and inverse.
      \item Reflect on why adding reasonable predictors should, in principle, reduce the RSE.
    \end{itemize}
  \end{alertblock}
\end{frame}

\end{document}
